\chapter{Homologia Morse'a: od punktów krytycznych do homologii singularnej}
\label{ch:morse-v10}

\section{Dlaczego rachunek punktów krytycznych wykrywa topologię?}
W teorii singularnej grupy łańcuchów są zwykle nieskończenie generowane,
ponieważ generatory stanowią wszystkie sympleksy singularne. Dla zwartej
gładkiej rozmaitości można skonstruować znacznie mniejszy kompleks:
jego generatorami będą punkty krytyczne jednej funkcji, a różniczkę
wyznaczymy, licząc odpowiednio zorientowane trajektorie gradientowe.
Zaskakujący jest fakt, że jego homologia nie zależy od wybranej
funkcji ani metryki: odzyskujemy zwykłą homologię singularną.
Przedstawimy dwa obrazy tej samej konstrukcji: geometryczny (trajektorie)
i komórkowy (doczepianie uchwytów).
Pełny geometryczny rachunek uchwytów znajduje się w
\href{https://github.com/dariusz-piekarz/mathematical_scipts/blob/biblioteka-matematyczna-iteracja-6/chapters/pl/differential_geometry_and_topology/18-teoria-morse-a-punkty-krytyczne-uchwyty-i-przeplyw.tex}{rozdziale \emph{Teoria Morse'a: punkty krytyczne, uchwyty i przepływ}}
skryptu z geometrii różniczkowej
(etykieta źródłowa \texttt{ch:teoria-morsea}).
Nie będziemy utożsamiać pełnego
twierdzenia o sklejaniu trajektorii z prostą konsekwencją twierdzenia
Sarda: wprowadza ono osobny, głęboki argument analityczny.

\section{Funkcja Morse'a, indeks i lokalny model}
\begin{definicja}[Punkt niezdegenerowany i indeks Morse'a]
Niech $M$ będzie gładką rozmaitością bez brzegu o wymiarze $n$,
a $f\in C^\infty(M,\R)$. Punkt $p\in M$ nazywamy \emph{krytycznym}, jeśli $\dd f_p=0$.
W dowolnym układzie współrzędnych wokół $p$ określmy w takim
punkcie formę $\operatorname{Hess}_p f$ przez macierz drugich
pochodnych funkcji zapisanej we współrzędnych.
Jest ona niezależna od wyboru współrzędnych jako
forma dwuliniowa na $T_pM$, gdyż
składniki zależne od pochodnych pierwszego rzędu znikają dla $\dd f_p=0$.
Punkt krytyczny jest \emph{niezdegenerowany}, gdy hesjan nie ma jądra.
\emph{Indeksem} $\lambda(p)$ nazywamy wymiar maksymalnej podprzestrzeni,
na której $\operatorname{Hess}_pf$ jest ujemnie określony.
Funkcja jest \emph{funkcją Morse'a}, jeżeli wszystkie jej punkty
krytyczne są niezdegenerowane.
Zbiór punktów krytycznych oznaczamy przez
$\operatorname{Crit}(f):=\{p\in M:\dd f_p=0\}$.
\end{definicja}

\begin{twierdzenie}[Lemat Morse'a — używane twierdzenie geometryczne]
\label{thm:morse-lemma-v10}
Jeśli $p$ jest niezdegenerowanym punktem krytycznym funkcji
$f\in C^\infty(M,\R)$ o indeksie $\lambda$, to istnieją gładkie
współrzędne $(x_1,\ldots,x_n)$, w których $p$ ma współrzędne $0$ oraz
\[
 f(x)=f(p)-\sum_{i=1}^{\lambda}x_i^2+
                \sum_{i=\lambda+1}^{n}x_i^2.
\]
W szczególności niezdegenerowane punkty krytyczne są izolowane.
\end{twierdzenie}
\begin{proof}
Przyjmijmy lokalne współrzędne z $p=0$. Ponieważ
$\mathrm df_0=0$, dwukrotne zastosowanie wzoru Newtona--Leibniza daje
\[
f(x)-f(0)=\int_{t\in[0,1]}(1-t)\,\D^2f(tx)[x,x] \,\dd t
=x^\top A(x)x,
\quad A(x)=\int_{t\in[0,1]}(1-t)\,\operatorname{Hess}f(tx)\,\dd t.
\]
Macierz $A(x)$ jest symetryczna i gładka oraz
$A(0)=\tfrac12\operatorname{Hess}_0f$ jest odwracalna.
Po jednej stałej liniowej zmianie współrzędnych
możemy założyć $A(0)=J=\operatorname{diag}(-I_\lambda,I_{n-\lambda})$.
Załóżmy najpierw $0<\lambda<n$ i rozpiszmy
$A(x)$ względem $\R^\lambda\oplus\R^{n-\lambda}$:
\[
 A(x)=\begin{pmatrix}P(x)&Q(x)\\ Q(x)^{\mathsf T}&C(x)\end{pmatrix}.
\]
W punkcie zerowym mamy $P(0)=-I_\lambda$,
$Q(0)=0$ oraz $C(0)=I_{n-\lambda}$.
Po zmniejszeniu otoczenia $P(x)$ pozostaje ujemnie
określona, a jej dopełnienie Schura
$S(x)=C(x)-Q(x)^{\mathsf T}P(x)^{-1}Q(x)$
pozostaje dodatnio określone. Wynika to z ciągłości
wartości własnych i równości $S(0)=I_{n-\lambda}$.
Rozpisanie iloczynu daje tożsamość
\[
 x^{\mathsf T}A(x)x=
 (x_-+P^{-1}Qx_+)^{\mathsf T}P(x_-+P^{-1}Qx_+)
 +x_+^{\mathsf T}Sx_+.
\]
Wyrazy mieszane mają bowiem sumę $2x_-^{\mathsf T}Qx_+$,
a kwadratowe wyrazy $Q^{\mathsf T}P^{-1}Q$ redukują się
z odpowiednim składnikiem $S$.
Zdefiniujmy zatem
\[
 y_-=(-P(x))^{1/2}(x_-+P(x)^{-1}Q(x)x_+),\qquad
 y_+=S(x)^{1/2}x_+.
\]
Dodatni pierwiastek kwadratowy macierzy dodatnio określonej
zależy gładko od jej współczynników: lokalnie blisko
$I$ ma rozwinięcie $\sqrt{I+E}=I+\tfrac12E-\tfrac18E^2+\cdots$,
które jest zbieżne wraz z pochodnymi na dostatecznie
małym otoczeniu. Ponieważ $P(0)=-I$, $Q(0)=0$,
$S(0)=I$, różniczka odwzorowania
$x\mapsto(y_-(x),y_+(x))$ w zerze jest identycznością;
pochodne macierzy są mnożone przez współrzędne $x$,
więc nie wnoszą dodatkowych składników. Z twierdzenia
o funkcji odwrotnej mamy lokalny dyfeomorfizm i postać
$f=f(0)-|y_-|^2+|y_+|^2$.
Przy $\lambda=0$ cały $A(x)$ jest dodatnio określony,
więc bierzemy $y=A(x)^{1/2}x$; dla $\lambda=n$
używamy $y=(-A(x))^{1/2}x$. W tych przypadkach
nie wprowadzamy bloków o zerowym wymiarze.
Ta postać pokazuje też, że $p$ jest izolowanym punktem krytycznym.
\end{proof}

\begin{wniosek}[Skończoność punktów krytycznych]
Jeżeli $M$ jest zwartą rozmaitością bez brzegu, a $f:M\to\R$ jest
funkcją Morse'a, to zbiór $\operatorname{Crit}(f)$ jest skończony.
\end{wniosek}
\begin{proof}
Zbiór krytyczny jest domknięty, gdyż jest miejscem zerowym ciągłej
sekcji $\dd f$ wiązki kostycznej, a zatem zwarty. Z poprzedniego
twierdzenia jest dyskretny. Zwarta przestrzeń dyskretna jest skończona.
\end{proof}

\section{Przepływ gradientowy, rozmaitości stabilne i transwersalność}
\begin{definicja}[Para Morse'a--Smale'a]
Niech $M$ będzie zwartą gładką rozmaitością bez brzegu,
$f:M\to\R$ funkcją Morse'a, a $g$ metryką riemannowską.
Pole $X=-\operatorname{grad}_g f$ ma zupełny przepływ
$\Phi:\R\times M\to M$. Dla $p\in\operatorname{Crit}(f)$ definiujemy
\begin{align*}
 W^u(p)&=\{x\in M:\lim_{t\to-\infty}\Phi_t(x)=p\},\\
 W^s(p)&=\{x\in M:\lim_{t\to+\infty}\Phi_t(x)=p\}.
\end{align*}
Para $(f,g)$ jest \emph{parą Morse'a--Smale'a}, jeżeli
$W^u(p)$ przecina $W^s(q)$ transwersalnie dla wszystkich
$p,q\in\operatorname{Crit}(f)$. Transwersalność rozumiemy w każdym
punkcie przecięcia, jako równość $T_xW^u(p)+T_xW^s(q)=T_xM$.
\end{definicja}

\begin{lemat}[Monotoniczność energii i wymiary]
\label{lem:morse-energy-v10}
Dla każdej trajektorii $t\mapsto\Phi_t(x)$ zachodzi
$\frac{\dd}{\dd t}f(\Phi_t(x))=-\|\operatorname{grad}_g f(\Phi_t(x))\|_g^2$.
Rozmaitości $W^u(p)$ i $W^s(p)$ są zanurzonymi rozmaitościami,
obie są dyfeomorficzne odpowiednio z $\R^{\lambda(p)}$ oraz
$\R^{n-\lambda(p)}$. Jeśli trajektoria nietrywialna łączy
$p$ z $q$, to $f(p)>f(q)$.
\end{lemat}
\begin{proof}
Pierwsza równość jest regułą łańcuchową i definicją gradientu.
Jeżeli całka z $\|\operatorname{grad}f\|^2$ na dowolnym odcinku
nietrywialnej trajektorii byłaby zerowa, gradient byłby na nim zerowy
i trajektoria byłaby stała, sprzeczność. Stąd ścisła nierówność.
W otoczeniu $p$ liniaryzacja pola ma postać
$-\operatorname{Hess}_pf$ po identyfikacji za pomocą $g$; ma
$\lambda(p)$ dodatnich i $n-\lambda(p)$ ujemnych wartości własnych.
Twierdzenie o rozmaitościach stabilnej i niestabilnej daje wskazane
wymiary i gładkie zanurzenia. Globalne rozmaitości otrzymuje się
przez obrazy lokalnych przy przepływie; w przypadku gradientowym
są dyfeomorficzne z przestrzeniami euklidesowymi odpowiednich wymiarów.
To ostatnie korzysta ze standardowego twierdzenia o rozmaitościach
stabilnych, które traktujemy jako narzędzie z geometrii dynamiki.
\end{proof}

Dla $p\ne q$ oznaczamy przez
$\widetilde{\mathcal M}(p,q)=W^u(p)\cap W^s(q)$
przestrzeń \emph{sparametryzowanych} trajektorii. Translacja czasu
$\tau\cdot x=\Phi_\tau(x)$ wyznacza wolne działanie $\R$.
Iloraz $\mathcal M(p,q)=\widetilde{\mathcal M}(p,q)/\R$
nazywamy przestrzenią \emph{trajektorii niesparametryzowanych}.

\begin{stwierdzenie}[Wymiar przestrzeni trajektorii]
\label{prop:morse-dimension-v10}
Dla pary Morse'a--Smale'a i $p\ne q$ przestrzeń
$\mathcal M(p,q)$ jest gładką rozmaitością wymiaru
$\lambda(p)-\lambda(q)-1$, o ile jest niepusta.
\end{stwierdzenie}
\begin{proof}
Transwersalność daje
$\dim(W^u(p)\cap W^s(q))=\lambda(p)+(n-\lambda(q))-n
=\lambda(p)-\lambda(q)$.
W każdym punkcie przecięcia pole $X$ jest niezerowe i styczne do
przecięcia, więc jego przepływ daje lokalny model
$(t,y)\mapsto\Phi_t(y)$ dla działania translacji.
Globalny przekrój tego działania otrzymujemy na przecięciu z
regularnym poziomem $f=c$ dla $f(q)<c<f(p)$:
na każdej trajektorii $f$ jest ściśle malejąca, więc przecina
poziom $c$ dokładnie raz, transwersalnie.
Ten przekrój jest dyfeomorficzny z ilorazem i ma wymiar
$\lambda(p)-\lambda(q)-1$.
\end{proof}

\section{Zliczanie trajektorii i kompleks Morse'a}
Wybierzmy orientację każdej niestabilnej rozmaitości $W^u(p)$.
Wzdłuż $W^s(q)$ wybór orientacji $W^u(q)$ indukuje orientację wiązki
normalnej do $W^s(q)$, zgodnie z rozkładem stabilnym i niestabilnym
przy $q$ i transportem przez przepływ. Gdy $\lambda(p)-\lambda(q)=1$,
w punkcie $x\in W^u(p)\cap W^s(q)$ transwersalność pozwala porównać
orientację $T_xW^u(p)$ z orientacjami normalną do $W^s(q)$
i dodatniego kierunku przepływu. Tak otrzymujemy znak
$\varepsilon(\gamma)\in\{-1,+1\}$ niesparametryzowanej
trajektorii $\gamma\in\mathcal M(p,q)$. To ujęcie nie wymaga
orientowalności całej rozmaitości $M$. Przy zmianie orientacji
$W^u(p)$ zmieniają się znaki trajektorii wychodzących z $p$;
przy zmianie orientacji $W^u(q)$ zmieniają się znaki trajektorii
dochodzących do $q$. Obie zmiany odpowiadają zmianie znaków
odpowiednich generatorów kompleksu Morse'a.

\begin{twierdzenie}[Kompaktyfikacja trajektorii — twierdzenie geometryczne]
\label{thm:morse-compact-v10}
Niech $M$ będzie zwartą gładką rozmaitością
bez brzegu, a $(f,g)$ parą Morse'a--Smale'a.
Dla dowolnych $p,q,r\in\operatorname{Crit}(f)$
spełniających warunki indeksowe poniżej zachodzi:
\begin{enumerate}[label=(\roman*)]
\item jeśli $\lambda(p)-\lambda(q)=1$, zbiór $\mathcal M(p,q)$
      jest skończony;
\item jeśli $\lambda(p)-\lambda(r)=2$, przestrzeń
      $\mathcal M(p,r)$ dopuszcza zwartą kompaktyfikację
      $\overline{\mathcal M}(p,r)$ o strukturze zorientowanej
      jednowymiarowej rozmaitości z brzegiem, przy czym
      \[
       \partial\overline{\mathcal M}(p,r)=
       \coprod_{\lambda(q)=\lambda(p)-1}
       \mathcal M(p,q)\times\mathcal M(q,r),
      \]
      a znak punktu brzegowego odpowiada iloczynowi znaków
      dwóch sklejanych trajektorii, z jednym wspólnym znakiem
      wynikającym z przyjętej orientacji brzegu.
\end{enumerate}
\end{twierdzenie}
\begin{proof}
Wyjaśnijmy dokładnie zakres użytego wyniku. Z równania energii
\[
\int_{t\in\R}\|\operatorname{grad}f(\gamma(t))\|_g^2\,\dd t
=f(p)-f(q)
\]
wynika ograniczenie energii trajektorii.
Przy zwartości $M$ ciąg trajektorii można, po dobraniu translacji
czasu i wybraniu podciągów, badać przez zwartość rozwiązań ODE
na przedziałach ograniczonych. Utrata zwartości może polegać na tym,
że trajektoria spędza coraz dłuższy czas w otoczeniu pośredniego
punktu krytycznego $q$. Powstaje wówczas trajektoria
\emph{łamana} $p\to q\to r$. Ścisły spadek $f$ i transwersalność ograniczają
liczbę łamań. Każdy niepusty zbiór trajektorii
nietrywialnych $\mathcal M(a,b)$ jest rozmaitością
wymiaru $\lambda(a)-\lambda(b)-1\ge0$,
więc każdy odcinek trajektorii łamanej obniża indeks
co najmniej o jeden. Przy całkowitej różnicy indeksów
równej $1$ nie może być dwóch takich odcinków;
przy różnicy $2$ jedyne możliwe łamanie ma
dokładnie jeden punkt pośredni $q$ i różnice
indeksów równe $1$ oraz $1$.
Aby stwierdzić, że każdemu takiemu złamaniu odpowiada
\emph{jeden koniec gładkiej jednowymiarowej rozmaitości}
$\mathcal M(p,r)$, potrzebne jest twierdzenie o sklejaniu
trajektorii i zgodności orientacji. Jego dowód wykorzystuje
lokalną postać przepływu przy punktach hiperbolicznych
oraz oszacowania dla sklejonych rozwiązań i nie jest
konsekwencją samej zwartości ODE. Przyjmujemy ten
konkretny wynik pomocniczy; wszystkie dalsze użycia
odnoszą się dokładnie do punktów (i)--(ii), bez
niejawnego założenia silniejszej kompaktyfikacji.
W szczególności nie twierdzimy, że samo ograniczenie energii
dowodzi istnienia współrzędnych sklejania lub zgodności
orientacji brzegowych. Te części są treścią przyjętego
twierdzenia geometrycznego.
\end{proof}

\begin{definicja}[Kompleks Morse'a]
\label{def:morse-complex-v10}
Dla pary Morse'a--Smale'a $(f,g)$ na zwartej rozmaitości bez brzegu
niech $CM_k(f,g)$ będzie wolną grupą abelową z bazą
$\{p\in\operatorname{Crit}(f):\lambda(p)=k\}$.
Określamy homomorfizm $\partial_k:CM_k(f,g)\to CM_{k-1}(f,g)$,
$\partial_kp=\sum_{\lambda(q)=k-1}n(p,q)q$, gdzie
$n(p,q)=\sum_{\gamma\in\mathcal M(p,q)}\varepsilon(\gamma)\in\Z$.
Skończoność sum wynika z Twierdzenia~\ref{thm:morse-compact-v10}.
\end{definicja}

\begin{twierdzenie}[Różniczka Morse'a spełnia $\partial^2=0$]
\label{thm:morse-d2-v10}
Dla każdej pary Morse'a--Smale'a $(f,g)$ na zwartej gładkiej
rozmaitości bez brzegu zachodzi $\partial_{k-1}\partial_k=0$
dla wszystkich $k\in\Z$.
\end{twierdzenie}
\begin{proof}
Współczynnik przy krytycznym punkcie $r$ o indeksie $k-2$
w $\partial_{k-1}\partial_kp$ wynosi
$\sum_{\lambda(q)=k-1}n(p,q)n(q,r)$.
Każdy składnik tego wyrażenia odpowiada punktowi brzegowemu
$\overline{\mathcal M}(p,r)$ z odpowiednim znakiem, według
Twierdzenia~\ref{thm:morse-compact-v10}(ii). Zorientowana
zwarta jednowymiarowa rozmaitość jest skończoną sumą odcinków
i okręgów; końce każdego odcinka mają przeciwne znaki,
a okręgi nie mają brzegu. Suma znaków jej punktów brzegowych
jest więc zerowa. Zatem wskazany współczynnik znika dla
każdego $r$, co dowodzi $\partial_{k-1}\partial_k=0$.
\end{proof}

Definiujemy $HM_k(f,g)=\ker\partial_k/\operatorname{im}\partial_{k+1}$.
Istotą teorii jest wykazanie, że grupy te nie są nowymi
niezmiennikami zależnymi od pary $(f,g)$.

\section{Od uchwytów do homologii singularnej}
\begin{twierdzenie}[Zmiana podpoziomu w punkcie krytycznym]
\label{thm:morse-handle-v10}
Niech $M$ będzie zwartą gładką rozmaitością bez brzegu,
$f:M\to\R$ funkcją Morse'a, a $a<b$ liczbami regularnymi
takimi, że $f^{-1}([a,b])$ zawiera dokładnie jeden punkt
krytyczny $p$ o indeksie $k$.
Wtedy podpoziom $M^b=f^{-1}(( -\infty,b])$ powstaje,
z dokładnością do homotopii, z $M^a$ przez dołączenie
jednej $k$-komórki. Ponadto
$H_j(M^b,M^a;\Z)\cong\Z$ dla $j=k$ i jest zerem
dla $j\ne k$.
\end{twierdzenie}
\begin{proof}
Z lematu Morse'a wokół $p$ mamy $f-f(p)=-\|u\|^2+\|v\|^2$
dla $u\in\R^k$, $v\in\R^{n-k}$. Uchwyt
$D^k\times D^{n-k}$ doczepiamy wzdłuż
$S^{k-1}\times D^{n-k}$, ponieważ tylko w tych kierunkach
występuje spadek wartości $f$ przy przechodzeniu
przez poziom krytyczny. Na obszarze bez punktów krytycznych, dla dowolnie ustalonej
metryki $g$, pole $-\operatorname{grad}_g f/\|\operatorname{grad}_g f\|_g^2$
spełnia $\dd f(X)=-1$; jego przepływ identyfikuje regularne
podpoziomy po pominięciu małych otoczeń punktów krytycznych.
Czas przepływu jest skończony, a zwartość $f^{-1}([a,b])$
zapobiega ucieczce trajektorii przez nieskończoność.
Daje to standardowe twierdzenie o dołączaniu uchwytu
(przyjmujemy jego lokalno-globalną wersję z teorii Morse'a).
Para $(D^k\times D^{n-k},S^{k-1}\times D^{n-k})$
deformacyjnie zwija się do $(D^k,S^{k-1})$.
Względna homologia tej pary jest $\Z$ w stopniu $k$
i zerem w pozostałych stopniach. Ekscyzja przenosi
ten rachunek na $(M^b,M^a)$.
\end{proof}

\begin{lemat}[Kilka punktów krytycznych o tej samej wartości]
\label{lem:morse-same-level-audit}
Niech $M$ będzie zwartą gładką rozmaitością bez brzegu,
$f:M\to\R$ funkcją Morse'a, a $a<c<b$ liczbami takimi, że
$a,b$ są regularne i jedyną wartością krytyczną w $[a,b]$ jest $c$.
Niech $p_1,\dots,p_r$ będą wszystkimi punktami krytycznymi na
poziomie $c$ i niech $k_i=\lambda(p_i)$. Wtedy $M^b$ powstaje z $M^a$,
z dokładnością do homotopii par, przez jednoczesne dołączenie
$r$ komórek wymiarów $k_1,\dots,k_r$. W szczególności
\[
 H_j(M^b,M^a;\Z)\cong
 \bigoplus_{\{i:k_i=j\}}\Z.
\]
\end{lemat}
\begin{proof}
Ponieważ punkty $p_i$ są skończenie liczne i izolowane, istnieją
parami rozłączne małe otoczenia Morse'a $U_i$ zawarte w
$f^{-1}((a,b))$. W każdym $U_i$ lokalny rachunek z dowodu
Twierdzenia~\ref{thm:morse-handle-v10} dostarcza uchwytu
$D^{k_i}\times D^{n-k_i}$ z częścią przyczepiającą
$S^{k_i-1}\times D^{n-k_i}$. Poza sumą tych otoczeń
nie ma punktów krytycznych na $f^{-1}([a,b])$;
znormalizowany przepływ gradientowy pozwala przesunąć
poziomy regularne do okolic $U_i$. Dlatego uchwyty można
doczepić jednocześnie, a ich części przyczepiające są
parami rozłączne. Względna ekscyzja redukuje
$H_j(M^b,M^a;\Z)$ do sumy prostej homologii par
$(D^{k_i}\times D^{n-k_i},S^{k_i-1}\times D^{n-k_i})$.
Każda wnosi jedną kopię $\Z$ w stopniu $k_i$.
Argument wykorzystuje tę samą geometryczną wersję twierdzenia
o doczepianiu uchwytów co Twierdzenie~\ref{thm:morse-handle-v10}.
\end{proof}

\begin{twierdzenie}[Zgodność homologii Morse'a i singularnej]
\label{thm:morse-singular-v10}
Dla zwartej gładkiej rozmaitości $M$ bez brzegu i pary
Morse'a--Smale'a $(f,g)$ istnieją izomorfizmy
$HM_k(f,g)\cong H_k(M;\Z)$ dla wszystkich $k\ge0$.
Są zgodne z odwzorowaniami kontynuacji pomiędzy
parami Morse'a--Smale'a, po ustaleniu zwyczajowych
identyfikacji z homologią singularną.
\end{twierdzenie}
\begin{proof}
Dla funkcji Morse'a wybierzmy przedziały rozdzielające
\emph{różne wartości krytyczne}, a niekoniecznie poszczególne
punkty krytyczne: kilka punktów może mieć tę samą wartość
(np. dwa siodła funkcji $\cos\theta+\cos\varphi$ na torusie).
Lemat~\ref{lem:morse-same-level-audit} umożliwia doczepienie
wszystkich uchwytów danego poziomu jednocześnie. Otrzymujemy
filtrację przez podpoziomy i strukturę CW, w której na każdy
punkt krytyczny $p$ przypada komórka wymiaru $\lambda(p)$.
Zatem grupa $k$-łańcuchów komórkowych jest wolna na punktach
krytycznych o indeksie $k$.
W transwersalnym ustawieniu Morse'a--Smale'a liczba
incydencji komórki niestabilnej $p$ z komórką
niższego indeksu $q$ jest stopniem odwzorowania przyłączającego,
a stopień ten można obliczyć przez zorientowane przecięcia
$W^u(p)$ z przekrojem poprzecznym do $W^s(q)$.
Każde takie przecięcie odpowiada jednej trajektorii
$p\to q$, a znak przecięcia jest dokładnie
$\varepsilon(\gamma)$. Otrzymujemy więc
$\partial^{\mathrm{cell}}p=\sum_qn(p,q)q=\partial p$.
W tym przejściu używa się \emph{twierdzenia o zgodności
liczb incydencji z liczeniem trajektorii}; wynika ono
z konstrukcji odwzorowań przyłączających za pomocą przepływu
oraz transwersalnego przecięcia z regularnym poziomem.
Jest to standardowe twierdzenie geometryczne, którego
pełny dowód wymaga kontroli deformacji brzegów
uchwytów w pobliżu wielu punktów krytycznych; nie
zastępujemy go bezpodstawnym stwierdzeniem, że
same grupy łańcuchów są z definicji identyczne.
Homologia komórkowa jest izomorficzna z singularną,
co kończy argument po zastosowaniu przytoczonego
twierdzenia o liczbach incydencji.
Zgodność tych identyfikacji z kontynuacją
jest standardową częścią twierdzenia o porównaniu
Morse'a z homologią singularną; jej mechanizm
opisujemy w następnej uwadze.
\end{proof}

\begin{uwaga}[Kontynuacja, niezmienniczość i zakres dowodu]
Dla dwóch par Morse'a--Smale'a $(f_0,g_0)$ i $(f_1,g_1)$
wybiera się gładką rodzinę danych i liczy trajektorie równania
$\dot\gamma=-\operatorname{grad}_{g_s}f_s(\gamma)$,
przy parametrze $s\in\R$ stałym dla $s\ll0$ i $s\gg0$.
Dla regularnej rodziny parametrów daje to odwzorowanie łańcuchowe
$CM_*(f_0,g_0)\to CM_*(f_1,g_1)$; kompaktyfikacja
jednowymiarowych przestrzeni trajektorii dowodzi zgodności
z różniczką. Porównując dwie homotopie rodzin,
otrzymuje się homotopię łańcuchową; odwrócenie
kontynuacji daje izomorfizm w homologii.
Pełne niezależne przeprowadzenie tej konstrukcji
wymaga analitycznego twierdzenia o sklejaniu
(zależnego od parametru). W naszym skrypcie
niezmienniczość wynika również z izomorfizmu
z homologią singularną w Twierdzeniu~\ref{thm:morse-singular-v10}.
\end{uwaga}

\section{Nierówności Morse'a i charakterystyka Eulera}
\begin{twierdzenie}[Silne nierówności Morse'a]
\label{thm:morse-inequalities-v10}
Niech $M$ będzie zwartą gładką rozmaitością
bez brzegu o wymiarze $n\ge0$, $(f,g)$
parą Morse'a--Smale'a, $m_k$ liczbą punktów krytycznych
indeksu $k$ i $b_k=\dim_\Q H_k(M;\Q)$.
Istnieje wielomian $Q(t)\in\Z[t]$ o nieujemnych współczynnikach taki, że
\[
\sum_{k=0}^{n}m_kt^k-\sum_{k=0}^{n}b_kt^k=(1+t)Q(t).
\]
W szczególności $m_k\ge b_k$,
$\sum_{j=0}^k(-1)^{k-j}m_j\ge
\sum_{j=0}^k(-1)^{k-j}b_j$ i
$\sum_k(-1)^km_k=\chi(M)$.
\end{twierdzenie}
\begin{proof}
Rozważmy kompleks łańcuchowy $CM_*(f,g)\otimes_{\Z}\Q$,
utworzony przez iloczyn tensorowy z $\Q$ nad $\Z$.
Definicję i własność uniwersalną przywołujemy z
\href{https://github.com/dariusz-piekarz/mathematical_scipts/blob/biblioteka-matematyczna-iteracja-6/chapters/pl/differential_geometry_and_topology/03-algebra-tensorowa-i-zewnetrzna.tex}{rozdziału o algebrze tensorowej}
(etykieta źródłowa \texttt{sec:iloczyn-tensorowy}).
Ponieważ $\Q=S^{-1}\Z$ dla $S=\Z\setminus\{0\}$,
a lokalizacja modułów jest funktorem dokładnym,
$\Q$ jest płaskim $\Z$-modułem. Stąd
$H_k(CM_*\otimes_{\Z}\Q)\cong HM_k(f,g)\otimes_{\Z}\Q
\cong H_k(M;\Z)\otimes_{\Z}\Q\cong H_k(M;\Q)$.
Nie utożsamiamy przy tym homologii wymiernej
z całkowitą: przejście do $\Q$ usuwa torsję.
Ponieważ $CM_k$ jest skończenie generowaną wolną grupą abelową,
jego wymiar po przejściu do $\Q$ wynosi $m_k$.
Niech
$r_k=\operatorname{rank}_\Q\partial_k$; przyjmujemy $r_0=r_{n+1}=0$.
Z wymiarów $\ker\partial_k$ i ilorazu homologii wynika
$m_k=b_k+r_k+r_{k+1}$. Pomnóżmy przez $t^k$
i zsumujmy po $k$. Otrzymujemy wzór z
$Q(t)=\sum_{k=0}^{n-1}r_{k+1}t^k$, którego współczynniki
są nieujemnymi liczbami całkowitymi.
Porównanie współczynników daje słabe nierówności.
Naprzemienne sumowanie do $k$ daje różnicę $r_{k+1}\ge0$;
podstawienie $t=-1$ daje wzór na charakterystykę Eulera.
\end{proof}

\section{Przykłady obliczeń}
\begin{przyklad}[Sfera i funkcja wysokości]
Dla $S^n\subset\R^{n+1}$, $n\ge1$, funkcja
$f(x)=x_{n+1}$ ma dwa punkty krytyczne: minimum indeksu $0$
i maksimum indeksu $n$. Dlatego $CM_0\cong CM_n\cong\Z$
i pozostałe grupy łańcuchowe są zerowe.
Dla $n>1$ różniczka automatycznie znika z przyczyn
stopniowania. Dla $n=1$ dwie trajektorie od maksimum
do minimum mają przeciwne znaki i znoszą się.
W obu przypadkach odzyskujemy $H_0(S^n;\Z)\cong\Z$
i $H_n(S^n;\Z)\cong\Z$.
\end{przyklad}

\begin{przyklad}[Cztery punkty krytyczne torusa]
Na $T^2=S^1\times S^1$ rozważmy
$f(\theta,\varphi)=\cos\theta+\cos\varphi$
z płaską metryką produktową. Punkty $(\pi,\pi)$,
$(0,\pi)$, $(\pi,0)$, $(0,0)$ mają indeksy odpowiednio
$0,1,1,2$. Stąd $CM_0\cong CM_2\cong\Z$ i $CM_1\cong\Z^2$.
Ponieważ znana homologia torusa ma rangi $1,2,1$,
nierówności Morse'a są tu równościami; wymusza to
$\partial_1=\partial_2=0$ po zastosowaniu iloczynu tensorowego z $\Q$,
a ponieważ grupy łańcuchowe są wolne, różniczki są zerowe
również nad $\Z$. Nie zakładaliśmy więc bez uzasadnienia,
że symetria automatycznie znosi wszystkie znaki.
\end{przyklad}

\begin{przyklad}[Niedoskonała funkcja na okręgu]
Dla $f(\theta)=\cos(2\theta)$ na $S^1$ mamy dwa maksima
$p_1,p_2$ i dwa minima $q_1,q_2$. Po wybraniu orientacji
niestabilnych kierunków różniczka może mieć postać
$\partial p_1=q_1-q_2$ oraz $\partial p_2=q_2-q_1$.
Macierz $\partial_1$ ma rząd $1$ nad $\Z$ i obraz
$\Z(q_1-q_2)$. Zatem $HM_1\cong\Z$ oraz
$HM_0\cong\Z$, mimo że funkcja posiada więcej punktów
krytycznych niż funkcja wysokości na okręgu.
\end{przyklad}

\begin{center}
\begin{tikzpicture}[>=Stealth,scale=0.93,font=\small]
\draw[->] (-3.55,0)--(3.58,0) node[right] {$\theta$};
\draw[very thick,manifoldblue,domain=-3.14159:3.14159,samples=140]
 plot (\x,{0.85*cos(deg(2*\x))});
\draw[manifoldred,->,thick] (-0.40,0.59)--(-0.75,0.06);
\draw[manifoldred,->,thick] ( 0.40,0.59)--( 0.75,0.06);
\fill[manifoldblue] (0,0.85) circle (1.7pt);
\fill[manifoldblue] (-1.5708,-0.85) circle (1.7pt);
\fill[manifoldblue] ( 1.5708,-0.85) circle (1.7pt);
\draw[manifoldblue,fill=white] (-3.14159,.85) circle (1.7pt);
\draw[manifoldblue,fill=white] ( 3.14159,.85) circle (1.7pt);
\node[above] at (0,1.02) {maksimum};
\node[below] at (-1.5708,-1.02) {minimum};
\node[below] at ( 1.5708,-1.02) {minimum};
\node[below] at (-3.14159,-0.08) {$-\pi$};
\node[below] at ( 3.14159,-0.08) {$\pi$};
\node[anchor=west] at (-3.46,1.53) {$f(\theta)=\cos(2\theta)$};
\end{tikzpicture}
\end{center}
\begin{center}\small Wykres $f(\theta)=\cos(2\theta)$ dla
$\theta\in[-\pi,\pi]$; końce przedziału oznaczają ten sam
punkt okręgu. Pełne punkty oznaczają trzy widoczne ekstrema
we wnętrzu przedziału, a punkty puste -- jeden wspólny
punkt maksymalny po utożsamieniu końców. Strzałki wskazują
kierunek ujemnego przepływu gradientowego.\end{center}

\begin{przyklad}[Przestrzeń rzutowa: homologia Morse'a wykrywa torsję]
\label{prz:morse-rp2-torsja}
Dla liczb rzeczywistych $\lambda_0<\lambda_1<\lambda_2$
określmy na $\R P^2$ funkcję
\[
 f([x_0:x_1:x_2])=
 \frac{\lambda_0x_0^2+\lambda_1x_1^2+\lambda_2x_2^2}
 {x_0^2+x_1^2+x_2^2}.
\]
Wzór nie zależy od reprezentanta, ponieważ licznik i mianownik
są jednorodne stopnia dwa. Wybierzmy metrykę ilorazową
z jednostkowej sfery $S^2$. Na sferze gradient ma składowe
$2(\lambda_i-f(x))x_i$, dlatego ujemny przepływ gradientowy
jest opisany jawnie przez
\[
 x_i(t)=\frac{e^{-2\lambda_i t}x_i(0)}
 {\left(\sum_{j=0}^2 e^{-4\lambda_jt}x_j(0)^2\right)^{1/2}}.
\]
Zera współrzędnych pozostają zerami, a niezerowe ilorazy
$x_i(t)/x_j(t)$ zmieniają się jak
$e^{-2(\lambda_i-\lambda_j)t}$.
Punkty krytyczne to dokładnie $p_i=[e_i]$, o indeksach
$\lambda(p_i)=i$: w karcie wokół $p_i$ hesjan ma znaki
$\lambda_j-\lambda_i$ dla $j\ne i$. Otrzymujemy
\[
 W^u(p_0)=\R P^0,\quad
 W^u(p_1)=\R P^1\setminus\R P^0,\quad
 W^u(p_2)=\R P^2\setminus\R P^1.
\]
Są to otwarte komórki $\R^i$; na ostatniej współrzędne
$x_0/x_2,x_1/x_2$ określają dyfeomorfizm z $\R^2$.
Powstaje więc dokładnie filtracja komórkowa
$\R P^0\subset\R P^1\subset\R P^2$.
Pierwsza komórka jednowymiarowa ma oba końce przyłączone
do jednego punktu, zatem $\partial_1=0$. Dysk
$D^2$ przyłączający górną komórkę ma brzeg, którego
punkty antypodyczne są utożsamione. Odwzorowanie
$S^1\to\R P^1\cong S^1$ ma stopień $\pm2$.

Sprawdźmy niezależnie znaki trajektorii. W karcie $x_2\ne0$
przyjmijmy $u=x_0/x_2$, $v=x_1/x_2$. Ujemne pole
gradientowe zapisuje się w tej karcie jako
\[
 \dot u=2(\lambda_2-\lambda_0)u,\qquad
 \dot v=2(\lambda_2-\lambda_1)v.
\]
Orientujemy $W^u(p_2)=\R^2$ bazą
$(\partial_u,\partial_v)$; stabilna rozmaitość $p_1$
przecina tę kartę w $\{u=0,v\ne0\}$.
Na poziomie regularnym $\lambda_1<c<\lambda_2$
dostajemy dokładnie dwa punkty przecięcia:
\[
 u=0,\qquad v=\pm\sqrt{\frac{\lambda_2-c}{c-\lambda_1}}.
\]
Odpowiadają one dwóm trajektoriom $p_2\to p_1$.
Na $W^u(p_1)$ wybieramy orientację wzrastającego
$a=x_0/x_1$. Ponieważ $a=u/v$, forma wyznaczająca
normalną do $W^s(p_1)$ ma znak $\dd a=\dd u/v$.
Dodatnio skierowany przepływ w $u=0$ ma znak
$v\partial_v$. Reguła orientacji „przepływ, potem
przekrój” ustala dodatni kierunek w przekroju
jako $-\partial_u$ dla $v>0$ i $+\partial_u$ dla $v<0$.
W obu punktach forma $\dd u/v$ ma na tak określonym
wektorze znak ujemny. Trajektorie mają zatem ten sam
znak, co daje $n(p_2,p_1)=-2$ dla przyjętej orientacji.
Po odwróceniu orientacji $W^u(p_2)$ otrzymujemy $+2$.

W karcie $x_1\ne0$ niestabilna rozmaitość $p_1$
to prosta $b=x_2/x_1=0$ ze współrzędną $a=x_0/x_1$.
Dwie półproste $a>0$ oraz $a<0$ wyznaczają dwie
trajektorie $p_1\to p_0$, ale kierunki przepływu
na nich są przeciwne względem orientacji wzrastającego
$a$. Ich liczby przecięć się znoszą:
$n(p_1,p_0)=0$. Zatem kompleks Morse'a ma postać
\[
 0\longrightarrow CM_2=\Z
 \xrightarrow{\;\partial_2=2\;} CM_1=\Z
 \xrightarrow{\;\partial_1=0\;} CM_0=\Z
 \longrightarrow0.
\]
Obliczamy kolejno $HM_2=\ker(2)=0$,
$HM_1=\ker(0)/2\Z\cong\Z/2\Z$ i $HM_0=\Z$.
Rachunek jest zgodny z komórkowym stopniem odwzorowania
przyłączającego, a jednocześnie stanowi osobne sprawdzenie
znaków trajektorii. Rangi grup łańcuchowych wynoszą
$(1,1,1)$, podczas gdy liczby Bettiego nad $\Q$
są równe $(1,0,0)$. Niezerowa różniczka o współczynniku
$2$ jest źródłem torsji w pierwszej homologii.
\end{przyklad}

\begin{figure}[!htbp]
\centering
\begin{tikzpicture}[>=Stealth,font=\small]
 \node[draw,rounded corners,fill=manifoldblue!8,
       minimum width=2.1cm,minimum height=1.15cm,
       align=center] (a) at (0,0) {$CM_2$\\$\Z$};
 \node[draw,rounded corners,fill=manifoldgreen!8,
       minimum width=2.1cm,minimum height=1.15cm,
       align=center] (b) at (4.0,0) {$CM_1$\\$\Z$};
 \node[draw,rounded corners,fill=manifoldgold!12,
       minimum width=2.1cm,minimum height=1.15cm,
       align=center] (c) at (8.0,0) {$CM_0$\\$\Z$};
 \draw[->,thick,manifoldred] (a)--(b)
       node[midway,above] {$\partial_2=2$};
 \draw[->,thick] (b)--(c)
       node[midway,above] {$\partial_1=0$};
 \node[below=6pt,align=center] at (4.0,-.6)
       {$HM_1=\operatorname{coker}(2:\Z\to\Z)=\Z/2\Z$};
\end{tikzpicture}
\caption{Kompleks Morse'a na $\R P^2$:
różniczka górnej komórki jest mnożeniem przez dwa,
a pierwsza homologia ma torsję.}
\label{fig:morse-rp2-torsja}
\end{figure}

\section{Związek z teorią uchwytów i z topologią algebraiczną}
Kompleks Morse'a pozwala obliczać niezmienniki topologiczne
przez skończony rachunek algebraiczny, a geometria przepływu
wyjaśnia, skąd biorą się współczynniki różniczki.
W rozdziale o Pontriaginie--Thomie punkty regularne
i podrozmaitości reprezentowały klasy stabilnej homotopii;
tutaj liczby trajektorii reprezentują homomorfizmy brzegowe w
kompleksie łańcuchowym. Przy późniejszym połączeniu z geometrią algebraiczną
pożyteczna będzie analogia z przestrzeniami modułów i
zorientowanymi liczbami przecięć, lecz do przeniesienia tej
konstrukcji potrzebne są osobne twierdzenia o zwartości,
transwersalności i orientacji przestrzeni modułów.


